\documentclass[12pt]{article}
\usepackage[utf8]{inputenc}
\usepackage[T1]{fontenc}
\usepackage{hyperref}
\title{Damal Language Resources for K-12 Curriculum Development}
\author{Compiled by Claude Sonnet 5 (claude-sonnet-5) for Andr\'as Kornai}
\date{}
\begin{document}
\maketitle

\section{Introduction}
Damal (also Amung, Uhunduni; ISO 639-3 uhn) is a Trans-New Guinea language of Papua, Indonesia, with about 14,000 speakers. Searches found little beyond Bible translation and audio materials. No school primers, textbooks or dictionaries were found.

\section{Primary Educational Resources}
No primers, textbooks or school programmes for Damal were found in search results.

\section{Linguistic and Reference Materials}
A general overview of the language is given at the Wikipedia article on the Uhunduni languages: \url{https://en.wikipedia.org/wiki/Uhunduni_languages}. No dictionary or grammar was confirmed in the searches.

\section{Religious and Community Resources}
The Christian and Missionary Alliance translation work began in 1957. The New Testament was printed in 1988 (\emph{Haik-A Ongam Kal}), and the whole Bible is reported complete (2021 MAF article): \url{https://maf-uk.org/?p=11637}. Scripture portal: \url{https://www.scriptureearth.org/00eng.php?iso=uhn}. Joshua Project profile: \url{https://joshuaproject.net/people_groups/11480}.

\section{Digital and Audio Resources}
Global Recordings Network offers Damal audio, including audio-visual Bible lessons: \url{https://globalrecordings.net/en/language/uhn}. A separate programme exists for the Enggipilu variety: \url{https://globalrecordings.net/en/language/9090}.

\section{Curriculum Development Context}
Damal is not known to be taught in schools. Any curriculum would need to begin with orthography and literacy materials, drawing on the Bible translation texts as the main body of written Damal. Little else is documented in the sources found.
\end{document}